\documentclass[11pt]{article}
\usepackage{explorer}
\newcommand{\ExplorerLanguage}{FR}
\newcommand{\ExplorerNumber}{06}
\newcommand{\ExplorerKind}{Exercices progressifs}
\begin{document}
\ExplorerTitle{Lycée · graphes et logique}{Parcours sans lever le crayon}
\ExplorerSection{Exercice 1}
Un graphe a des degrés 2,3,3,4. Combien d’arêtes possède-t-il ?\par
\vspace{5pt}
\ExplorerSection{Exercice 2}
Prouver que le nombre de sommets de degré impair est pair.\par
\vspace{5pt}
\ExplorerSection{Exercice 3}
Peut-on tracer un carré et ses deux diagonales sans repasser sur une arête ? Les intersections intérieures ne sont pas des sommets.\par
\vspace{5pt}
\ExplorerSection{Exercice 4}
Pourquoi deux cycles disjoints ne permettent-ils pas un parcours unique, alors que tous les degrés sont pairs ?\par
\vspace{5pt}
\end{document}
